\section*{Discussion}

The evaluation should be read as a test of a framework prototype for comparing VR prism-effect implementations. From that perspective, the main finding is that the sandbox made several normally separate parts of VR prism-adaptation work comparable inside one workflow: perturbation type, input mode, task, exposure quality, response confirmation, and baseline-to-post analysis. This supports the central motivation of the project. A shared framework is useful because the term ``VR prism effect'' is too broad on its own; studies also need to describe what is shifted, when feedback is visible, how responses are confirmed, and how after-effects are measured.

The non-zero None condition is especially important for this framework argument. It shows why a sandbox should include a no-effect reference rather than only plotting perturbation magnitudes. Baseline-to-post shifts can occur without an active perturbation, likely due to ordinary drift, fatigue, practice, calibration, or carry-over. The participant-specific normalization in Figure~\ref{fig:article-results-normalized} is therefore more informative than absolute magnitude alone. Even so, the small sample size limits the strength of any condition-level conclusion. Ten participants generated many logs, but only five participants contributed to each input-mode group.

The embodied condition also showed why implementation details belong in the same framework as the experimental task. Hand tracking was usable in some cases, but it introduced variability through tracking stability and hand posture. This is consistent with Cho et al., who noted practical hand-tracking limitations such as occlusion and self-occlusion in a VR prism-adaptation setup \cite{cho2022feasibility}. In the present prototype, small differences in hand posture could produce large changes in projected aim direction. Controller input was more robust, but controller triggers can disturb the endpoint at confirmation, which is why the final design used opposite-hand confirmation.

The movement itself also differs from conventional physical prism-adaptation procedures. Conventional procedures often use physical pointing or reaching from the body midline toward a target \cite{serino2006mechanisms,Serino2011}. The prototype used a raycast from the virtual controller or hand representation. This made the task feasible and consistent across wall-based targets, but it allowed participants to aim through wrist rotation, controller orientation, arm posture, or upper-body movement. The system therefore measured a spatial response shift, but the underlying motor process may not be identical to reaching-based prism adaptation. This is not only a limitation of the prototype; it is also an example of the terminology problem the framework is meant to expose.

The exposure procedure used concurrent feedback and counted misses as attempts. Terminal and concurrent feedback have been shown to produce different outcomes in prism-adaptation procedures \cite{Serino2011}, and Cho et al. implemented a VR variant where the hand was visible only in a target-visible area \cite{cho2022feasibility}. The present evaluation used concurrent feedback to keep the visible perturbations understandable and comparable across the implemented effects. This choice should be reported explicitly because feedback visibility changes what kind of adaptation procedure is being implemented. Similarly, 60 exposure attempts should not be interpreted as 60 successful corrections when hit rates vary. This exposure dose was also shorter than several reviewed protocols, including 96 reaches in Anan et al. and 100 reaches in Carter et al. and Bourgeois et al. \cite{anan2025comparative,Carter2016,Bourgeois2021}.

Session tolerance was another procedural constraint. The SSQ did not indicate severe simulator sickness as the main evaluation problem, but repeated raised-arm interaction produced fatigue and reduced focus in some participants. Future studies should therefore treat sickness and physical fatigue as separate measurements rather than assuming one explains the other.

The evaluation was also limited to one session per participant. This is appropriate for testing whether the sandbox can compare immediate after-effect data across prism-effect implementations, but it does not address long-term retention or repeated-session change. Those questions belong to later studies that use the framework in a rehabilitation-oriented protocol.
